\section*{\centering \fontsize{35}{45}\selectfont \fontsize{35}{45}\selectfont Basic Rules}
\vspace*{0.4 cm}

\begin{multicols*}{2}
\label{sec:basic}

\subsection*{Focus \& Stress}

There are two important resources each knight has:
the amount of thought they can put into perfecting a task and the amount of effort they can put into a failed task to fix it.
Those are tracked by \texttt{Focus} and \texttt{Stress}, respectively.

\paragraph{Focus.} Before a roll is performed, if it can be affected by a character, they can spend Focus to
enhance it. It \textbf{must} be declared \textbf{before} rolling the dice. For each point of \texttt{Focus} spent -- add \texttt{Advantage} to that roll.

One can not only plot how to enhance the roll, but also how to sabotage a roll.
If one character can affect the roll of another destructively by planning in advance, for each \texttt{Focus} spent, the roll accumulates \texttt{Disadvantage}.
\emph{Defense} governs this interaction for the attack rolls(p. \pageref*{sec:combat}).

Usually, a knight has no restriction on how much \texttt{Focus} they can spend to enhance an action, but
determining how much focus is possible to spend to oppose a given action should be discussed with the Referee. 

\paragraph{Stress.} After a roll is made, a knight can try to improve it by acquiring \texttt{Stress}.
For each gained \texttt{Stress}, add \texttt{Advantage} to that roll. Increasing one's \texttt{Stress} \textbf{must} be done \textbf{after} the roll.

Akin to \texttt{Focus}, one can also acquire \texttt{Stress} to oppose an action after the roll is made.

Dealing with an explicit incoming danger -- any amount of \texttt{Stress} can be spent.
For a split-second reaction, like avoiding a sudden fall or dodging an arrow -- up to \texttt{1} \texttt{Stress} can be spent.

A character can't spend \texttt{Stress} if they are at full \texttt{Stress}. 

\colbox{
    \paragraph{Mental breakdown} happens when \texttt{Stress} is full after roll is resolved.

    \begin{enumerate}[nosep, leftmargin=*]
        \item Reduce \texttt{Stress} by \textbf{1}
        \item Lose \texttt{d6} \texttt{Spirit}
        \item Suffer \textbf{Trauma}(p.\pageref*{sec:injuries})
    \end{enumerate}
}
If the action is contested by \texttt{Focus} or \texttt{Stress}, opposing characters reveal how much they spend simultaneously for each. 

\newcolumn

\subsection*{Fear \& Hope}

\label{sec:hopefear}
The battle rages as much on the field as in the warrior's mind. 
Extinguishing fear and tending to dwindling hope is as important as finding weaknesses in the enemy's armour. 

\paragraph{Tracking NPC state.} Usually, non-player characters(NPCs), are tracked with less granularity than those controlled by players. 
Main stats that are not tracked are \texttt{Injuries} and \texttt{Stress}: for bookkeeping and get-out-of-jail nature, respectively. 

Instead of tracking \texttt{Injuries}, each time an NPC takes \texttt{Vigor} damage that is half, rounding up, of their remaining \texttt{Vigor} -- they blackout. 
They die if left untreated for a while.

Instead of tracking \texttt{Stress} of hostile NPCs Referee uses \texttt{Fear}. The players can use \texttt{Hope} to do the same for characters' allies.

\paragraph{Fear.} Each time a character or a friendly npc rolls \texttt{1} on their weapon die or \texttt{20} on a Virtue Save the Company accumulates one point of \texttt{Fear}. 
This is a resource for the Referee, not characters.

When a hostile NPC wants to use \texttt{Stress} -- the Referee can spend \texttt{Fear} of the Company instead of the NPC's \texttt{Stress}. 
NPCs don't acquire \texttt{Trauma}s; how much it is suitable to spend is up to the Referee.

Additionally, if a character attempts something risky -- it might make sense narratively to manifest their \texttt{Fear} as if they were opposed by some external forces. 

Fear is a property of the Company, so it is persistent across passage of time and scenes.

\paragraph{Hope.} Each time an enemy of the Company rolls \texttt{1} on their weapon die or \texttt{20} on a Virtue Save the Company accumulates one point of \texttt{Hope}.

This is a resource, but this time -- for the Company. 
Any player can use \texttt{Hope} for \emph{Defense}, as \texttt{Focus} for themselves or as \texttt{Stress} for friendly NPCs. 
Spending \texttt{Hope} can't cause \texttt{Trauma}.

Hope is also a property of the Company, so it also persists. 

\end{multicols*}
